%==============================================================================
% Capítulo 1: Introducción
%==============================================================================
\section{Introducción}\label{sec:introduccion}

\subsection{Contextualización}\label{subsec:contextualizacion}

La evolución de las tecnologías de la información ha transformado profundamente la dinámica operativa de las organizaciones contemporáneas. En la actualidad los sistemas informáticos constituyen el soporte fundamental para la ejecución de procesos administrativos, productivos y de comunicación, y esa dependencia tecnológica ha incrementado de manera correlativa la importancia de los servicios de soporte técnico, particularmente aquellos relacionados con la gestión de incidentes informáticos. La disciplina que da marco a esa actividad ---la gestión de servicios de tecnologías de la información (ITSM)--- se desarrolla en la \autoref{sec:marco-teorico}, donde se define su alcance y su relación con la mesa de ayuda.

No obstante, en muchas organizaciones medianas el proceso de registro inicial de incidentes continua realizándose manualmente, lo que introduce demoras, inconsistencias y una marcada dependencia del criterio individual del operador.

El contexto argentino presenta características particulares que acentúan esta problemática. Las organizaciones medianas del país ---segmento que la Secretaría de la Pequeña y Mediana Empresa encuadra por tramos según el sector y parámetros como el personal ocupado o el volumen de ventas anuales \parencite{SEPyME2026_mipyme}--- operan con presupuestos de tecnología acotados y enfrentan restricciones cambiarias que encarecen la adopción de soluciones comerciales internacionales. La combinación de estas restricciones con una demanda creciente de servicios digitales configura un escenario en el que las soluciones basadas en componentes de código abierto y servicios de inferencia bajo demanda resultan particularmente atractivas desde el punto de vista económico y operativo.

\subsection{Planteamiento del problema}\label{subsec:planteamiento-problema}

La presente investigación se desarrolla en el marco de una organización mediana del sector servicios con sede en la provincia de Mendoza, República Argentina, cuya plantilla asciende a aproximadamente 120 usuarios internos distribuidos en cinco áreas funcionales. El área de operaciones desempeña simultáneamente tareas vinculadas a la gestión operativa cotidiana y a la atención de incidentes informáticos reportados por los usuarios, lo que obliga al personal responsable a interrumpir actividades críticas para registrar manualmente fallas de hardware, inconvenientes de software o solicitudes de asistencia técnica.

Las mediciones preliminares realizadas durante una semana laboral representativa evidenciaron un volumen promedio de 42 incidentes diarios, un tiempo medio de registro manual de 2 minutos con 45 segundos por incidente y una tasa de derivación erronea inicial cercana al 15~\%. Aquel valor de 2 minutos con 45 segundos constituyó una estimación a priori, obtenida del promedio de una semana completa de operación sin cronometraje individual por caso, y sirvió para dimensionar el problema. La campaña experimental posterior midió el tiempo de registro incidente por incidente con cronometraje en vivo, consignado en la planilla descrita en la \autoref{subsec:diseno-experimental}, y arrojó una media de 61,6~s sobre los 200 casos del corpus, sustancialmente inferior a la estimación preliminar. Sobre estos indicadores, el tiempo que el personal podría redirigir a tareas de mayor valor agregado si la etapa de registro estuviera automatizada se estima en el orden de 32,2 horas-persona por trimestre (véase \autoref{subsec:interpretacion-resultados}). Este tiempo representa un costo de oportunidad significativo para la organización.

El registro manual de incidentes implica además la lectura de correos electrónicos, la interpretación de descripciones textuales heterogeneas, la clasificación del problema según el área responsable y la carga manual de la información en un sistema de tickets. Este procedimiento requiere tiempo, interrumpe las tareas del personal técnico y puede generar errores de clasificación que prolongan la cadena de resolución cuando un ticket es derivado a un área incorrecta. La dispersión de los canales de comunicación ---correo, llamada telefónica y formularios internos--- agrava la situación al obligar al personal a monitorear multiples fuentes simultáneamente.

Este escenario responde a un patrón documentado en la literatura: \textcite{Galup2009_itsm} advierten que los procesos de soporte no automatizados tienden a degradarse cuando crece el volumen de solicitudes, con un impacto desproporcionado en organizaciones medianas donde el personal combina multiples roles (véase \autoref{subsec:gestion-servicios-ti}). Estos indicadores justifican la búsqueda de una alternativa tecnológica que reduzca el costo operativo, mejore la calidad de la derivación temprana y disminuya la variabilidad asociada al criterio individual.

\subsection{Pregunta de investigación}\label{subsec:pregunta-investigacion}

A partir del problema descrito, esta investigación se orienta a responder la siguiente pregunta: ?`reduce la automatización del registro inicial de incidentes mediante orquestación de flujos y procesamiento de lenguaje natural el tiempo de registro respecto del proceso manual, y alcanza la exactitud de clasificación el umbral objetivo definido para una mesa de ayuda empresarial mediana?

La pregunta se descompone en tres interrogantes subsidiarios que guian el diseño experimental y la selección de métricas. Primero, ?`cual es la magnitud de la reducción del tiempo de registro que puede obtenerse mediante la automatización, y resulta esa reducción estadísticamente significativa? Segundo, ?`alcanza el clasificador automático el umbral objetivo de exactitud global del 85~\% y un valor F1 macro no inferior a 0,85, y como se distribuyen los errores entre las categorías objetivo? Tercero, ?`permite la arquitectura propuesta integrar canales multiples de recepción sin degradar la latencia ni la calidad de la clasificación?

\subsection{Hipótesis de trabajo}\label{subsec:hipotesis-trabajo}

La hipótesis principal de este trabajo se desdobla en tres sub-hipótesis contrastables, cada una con su propio contraste y su criterio de falsación, y no se enuncia como una conjunción indivisible de mecanismo y umbral de desempeño. La primera sostiene que la automatización reduce el tiempo de registro: su contraste es el no paramétrico de Wilcoxon sobre las medianas, con \(H_{0}: M_{\text{manual}} = M_{\text{automatizado}}\) frente a \(H_{1}: M_{\text{manual}} > M_{\text{automatizado}}\), donde \(M\) representa la mediana poblacional del tiempo de registro, y su criterio de falsación es no rechazar \(H_{0}\). El contraste se rechaza con los valores vigentes \(W = 363\), \(p = 3{,}08\text{e-}32\) y tamaño del efecto \(r = 0{,}9639\) (\autoref{sec:resultados}); esta sub-hipótesis se declara confirmada.

La segunda postula que la automatización no incrementa la proporción de intervención humana respecto del proceso manual: su contraste previsto es la proporción de casos derivados a revisión por gate contra la fracción esperada, y su criterio de falsación es que esa proporción no resulte inferior a la esperada. No se realiza con la evidencia disponible ---no existe un registro primario de intervención por caso y el gate de confianza (inferior a 0,70) no deriva casos en el corpus---, de modo que se declara \emph{no contrastada por ausencia de registro primario} y no se cuantifica (véase \autoref{subsec:limitaciones-estudio}).

La tercera postula que la exactitud y el F1 macro alcanzan el umbral de desempeño (exactitud no inferior al 85~\% y F1 macro no inferior a 0,85), verificado contra la exactitud observada y su intervalo de confianza de Wilson, sin invocar la zona de rechazo del contraste de tiempos; su criterio de falsación es una exactitud observada por debajo del 85~\% o un intervalo de confianza que no supera el umbral. El resultado vigente ---73,5~\% de exactitud, IC \([67{,}0~\%, 79{,}1~\%]\) y macro-F1 de cinco clases de 0,5207 (\autoref{sec:resultados})--- no alcanza el umbral; esta sub-hipótesis se declara refutada.

Como hipótesis subsidiarias se plantea, en primer lugar, que la integración de multiples canales de entrada dentro de un único flujo automatizado mejora la accesibilidad del sistema y reduce la dispersión operativa sin introducir cuellos de botella adicionales. En segundo lugar, que un esquema híbrido de clasificación basado en reglas determinísticas y modelo de lenguaje grande resulta más eficiente ---en términos de la relación entre exactitud, latencia y costo de inferencia--- que un esquema construido exclusivamente sobre modelo externo, dado que una proporción sustancial de los incidentes presenta marcadores léxicos inequivocos que no requieren consulta al modelo de lenguaje.

\subsection{Objetivo general}\label{subsec:objetivo-general}

Evaluar la viabilidad y el desempeño de la automatización del registro inicial de incidentes de una mesa de ayuda empresarial mediana mediante la orquestación de flujos con N8N y el procesamiento de lenguaje natural en Python. El desarrollo y la implementación del sistema constituyen el medio para esa evaluación, que abarca la recepción, el procesamiento, la clasificación y el registro automáticos de incidentes y considera la reducción del tiempo de registro, la calidad de la clasificación y la proporción de intervención humana en la etapa inicial, manteniendo la supervisión técnica y la trazabilidad de las decisiones tomadas por el sistema.

\subsection{Objetivos específicos}\label{subsec:objetivos-especificos}

El primer objetivo específico consiste en diseñar una arquitectura distribuida y modular que integre un motor de orquestación de flujos, un módulo de procesamiento de lenguaje natural, un canal de telefonía con transcripción automática del habla y una base de datos relacional, garantizando comunicación asincrónica entre componentes mediante interfaces de programación de aplicaciones de tipo REST sobre HTTP cifrado.

El segundo objetivo específico se orienta a implementar un clasificador automático que asigne cada incidente a uno de los cinco sectores canónicos ---Seguridad Informática, Soporte Técnico Hardware, Soporte Técnico Software, Bases de Datos o Sistemas--- sobre un corpus controlado de 200 casos representativos del entorno operativo. El umbral de desempeño ---exactitud global no inferior al 85~\% y F1 macro no inferior a 0,85--- no es parte definitoria de este objetivo, sino un criterio de éxito contingente verificable contra la exactitud observada y su intervalo de confianza; su grado de cumplimiento se declara en la \autoref{subsec:cumplimiento-objetivos}. Ese umbral se fundamenta en tres criterios convergentes. Primero, los trabajos de \textcite{Paramesh2019_it_service_desk} reportan exactitudes de aproximadamente 87~\% con enfoques de máquinas de vectores de soporte (SVM) y representaciones TF-IDF, estableciendo un piso de referencia para el estado del arte en clasificación de tickets. Segundo, el costo organizacional de una derivación erronea se estima en el orden de 10 a 15 minutos adicionales de reasignación, estimación propia de los autores a partir de la observación operativa de la organización adoptante, sin un registro cuantitativo conservado, por lo que una exactitud inferior al 85~\% resultaría en tiempos de resolución comparables a los del flujo manual para un volumen de errores inaceptable. Tercero, situar el umbral en el 85~\% deja margen para que, si la exactitud verdadera del clasificador lo supera, el límite inferior de su intervalo de confianza pueda ubicarse por encima de la frontera de interés en lugar de quedar solapado con ella. Se trata, en consecuencia, de un criterio de desempeño verificable contra la exactitud observada y su intervalo de confianza, distinto del contraste sobre la mediana de los tiempos formulado en la \autoref{subsec:hipotesis-trabajo}.

El tercer objetivo específico plantea integrar tres canales de entrada paralelos ---correo electrónico, formulario web y llamada telefónica con transcripción automática--- dentro de un único flujo automatizado que normalice la entrada y centralice el procesamiento, mejorando la accesibilidad del sistema sin incrementar la complejidad operativa percibida por los usuarios finales.

El cuarto objetivo específico consiste en evaluar comparativamente el flujo manual y el flujo automatizado en términos de tiempo medio de registro, dispersión, exactitud de clasificación y proporción de intervención humana, mediante un diseño cuasiexperimental con muestras pareadas, prueba estadística no paramétrica y reporte de intervalos de confianza al 95~\%. El diseño experimental sigue las recomendaciones de \textcite{HernandezSampieri2018_metodologia} para investigaciones tecnológicas de carácter aplicado.

El quinto objetivo específico se orienta a documentar la arquitectura, los procedimientos de despliegue y los resultados experimentales de modo que una tercera parte pueda replicarlos, y a proponer una hoja de ruta de evolución que oriente futuras mejoras y replicaciones del estudio en organizaciones de características comparables.

Como sexto objetivo específico se plantea analizar los errores de clasificación y calibrar la confianza de las predicciones, tipificando los patrones de error y caracterizando las bandas de confianza para interpretar el desempeño más allá de las métricas agregadas. Su cumplimiento es verificable contra el reporte de evaluación del repositorio y la \autoref{subsec:analisis-errores} de la \autoref{sec:resultados}.

\subsection{Justificación}\label{subsec:justificacion}

La justificación del trabajo se sostiene en tres dimensiones complementarias.

Desde la dimensión organizacional, la automatización propuesta libera capacidad operativa del personal técnico, reduce el tiempo de respuesta percibido por los usuarios y disminuye la variabilidad introducida por el criterio individual. La evidencia recogida en este trabajo permite cuantificar esa mejora en el contexto específico de una organización mediana argentina, aportando un caso documentado que otras organizaciones pueden utilizar como referencia para estimar el retorno esperado de una inversión similar.

Desde la dimensión tecnológica, la solución combina componentes de código abierto ---Python, FastAPI y PostgreSQL--- con N8N, una plataforma de orquestación de flujos cuyo código fuente se distribuye bajo licencia Sustainable Use (\emph{fair-code}) y admite despliegue completamente autoalojado \parencite{n8n2024_docs}, lo que permite una implementación reproducible en organizaciones de tamaño mediano sin requerir grandes inversiones de capital ni dependencia exclusiva de proveedores comerciales. Esta propiedad es particularmente relevante en el contexto argentino, donde las restricciones de acceso a divisas y los costos de suscripción en moneda extranjera constituyen barreras significativas para la adopción de soluciones comerciales internacionales. Además, el despliegue autoalojado de la totalidad de los componentes ---con excepción del servicio de inferencia--- otorga a la organización control pleno sobre sus datos operativos, en línea con el control sobre los datos personales y su tratamiento conforme a la normativa de protección de datos \parencite{Ley25326_2000}.

Desde la dimensión académica, el trabajo aporta evidencia empírica sobre el desempeño de modelos de lenguaje grandes aplicados a la clasificación de tickets en español rioplatense, dominio escasamente documentado en la literatura comparativa disponible. Buena parte de la literatura relevada sobre clasificación automática de tickets se apoya en corpus en inglés y, en menor medida, en alemán o portugués \parencite{zangari2023ticket,colavito2025benchmarking}, y los pocos trabajos que incorporan el español lo hacen con variantes peninsulares o con corpus genéricos no específicos del dominio de mesas de ayuda \parencite{grandury2025laleaderboard,mayorrocher2026spanish}. Este trabajo contribuye a cerrar esa brecha mediante la validación rigurosa de un clasificador híbrido sobre un corpus de 200 casos en español rioplatense con doble etiquetado independiente y acuerdo inter-evaluador cuantificado.

El relevamiento propio de la documentación pública de los proveedores (consultada el 19-10-2026) muestra que plataformas comerciales como Zendesk Suite, Freshdesk o ServiceNow ITSM operan como software como servicio, con datos en jurisdicciones extranjeras y costos por agente y mes que escalan con el equipo \parencite{Zendesk2026_help_desk,Freshdesk2026_platform,ServiceNow2026_itsm}. La propuesta desarrollada en esta investigación utiliza herramientas de despliegue autoalojado, lo que facilita su adopción en organizaciones medianas y simplifica el cumplimiento de la normativa argentina de protección de datos personales.

\subsection{Alcance y delimitaciones}\label{subsec:alcance-delimitaciones}

El alcance del trabajo se circunscribe a la etapa de recepción, clasificación y registro inicial de incidentes, sin abarcar la resolución técnica de los mismos, la cual permanece bajo responsabilidad del personal humano. El sistema contempla tres canales de entrada (correo electrónico, formulario web y llamada telefónica con transcripción automática) y cinco sectores de derivación (Seguridad Informática, Soporte Técnico Hardware, Soporte Técnico Software, Bases de Datos y Sistemas). La validación experimental se realiza sobre un corpus en idioma español rioplatense recolectado en una organización mediana del sector servicios con sede en la provincia de Mendoza, de aproximadamente 120 usuarios internos distribuidos en cinco áreas funcionales.

Quedan fuera del alcance los siguientes aspectos, los cuales se proponen como líneas de trabajo futuro en la \autoref{sec:recomendaciones}. Primero, la integración con sistemas externos de inventario de activos informáticos o con sistemas de gestión de configuración. Segundo, los modelos de aprendizaje supervisado entrenados con corpus propios de la organización, dado que el volumen actual de datos etiquetados resulta insuficiente para esa estrategia. Tercero, la implementación de paneles de monitoreo avanzados con visualización en tiempo real de métricas operativas. Cuarto, la automatización de la etapa de resolución técnica, que trasciende el alcance del presente trabajo.

Las delimitaciones del estudio incluyen: la restricción geográfica a una única organización de la provincia de Mendoza, lo que limita la validez externa de los hallazgos; la dependencia operativa del servicio externo de inferencia gemini-3.6-flash, cuyo rendimiento y condiciones comerciales escapan al control de los investigadores; y la naturaleza temporal de los resultados, que reflejan el comportamiento del modelo en la versión vigente al momento del estudio (entre junio y octubre de 2026), sujeta a modificaciones futuras por parte del proveedor.

\newpage
